\section{관련 연구}
\label{sec:related-work}

\subsection{최종 순위 예측에서 페이즈별 생존 예측으로}

PUBG 예측 연구의 출발점은 경기 종료 후의 최종 순위였다. Kaggle의 PUBG Finish
Placement Prediction 데이터셋은 처치 수, 피해량, 이동 거리, 회복 아이템 사용량,
매치 유형 같은 경기 후 집계값으로 \texttt{winPlacePerc}를 예측한다
\cite{kagglePubgFinishPlacement}. Rokad 등, Sen 등, Lee와 Lee의 연구도 PUBG의
최종 순위와 관련된 변수와 예측 가능성을 다루었다
\cite{rokad2019pubgPlacement,sen2021pubgFinalRank,lee2025pubgRank}. 이 연구들은
PUBG 순위가 로그 기반 변수로 예측 가능하다는 점을 보여주지만, 대부분의 접근은
매치를 이미 끝난 하나의 기록으로 압축한다. 페이즈별 생존 문제에서는 이 압축이
강한 단순화가 된다. 같은 총 피해량이나 이동 거리도 전혀 다른 노출 이력에서 나올
수 있기 때문이다.

Sagong 등의 연구는 이러한 압축을 풀기 시작한 중요한 선행 작업이다
\cite{sagong2025hunterCamper}. 이들은 PUBG 로그를 페이즈별로 나누고, 공격성과
활동성을 나타내는 점수로 Hunter, Sniper, Traveler, Camper 같은 전략 유형을
구성한 뒤 다음 페이즈 생존과 승리를 예측했다. 이 논문은 그 문제의식을 이어받되,
전략 유형보다 앞선 표현 단위로 돌아간다. 팀의 상태, 이동, 안전구역 맥락을 어떤
절단 시점에서 관측해야 다음 페이즈의 탈락 위험 문제가 되는지를 묻는다.

\subsection{텔레메트리, 이동 이력, 동적 공간 맥락}

게임 텔레메트리 연구는 행동 로그가 단순 이벤트 기록을 넘어 플레이어 행동을
정량화하는 자료가 될 수 있음을 보여준다. Drachen 등은 행동 로그 기반 군집화와
프로파일 해석을 다루었고 \cite{drachen2014behavioralTelemetry}, Bunian 등은 행동의
순서가 집계 횟수보다 플레이어 차이를 더 잘 드러낼 수 있음을 보였다
\cite{bunian2018hmmGameBehavior}. 배틀로얄 장르의 게임에서는 이 순서성이 특히
중요하다. 예측 시점 이후의 피해량, 위치, 탈락 사건이 입력에 섞이면 모델은 생존을
예측하는 것이 아니라 이미 일어난 결과를 확인하게 된다.

이동 궤적 분석은 움직이는 대상을 시간순 공간 표본으로 다루는 언어를 제공한다
\cite{wang2020trajectorySurvey}. 팀 스포츠의 시공간 분석도 개별 선수의 행 단위
통계만으로는 결과를 충분히 설명하기 어렵고, 위치와 상호작용을 함께 보아야 함을
보여준다 \cite{gudmundsson2016teamSportsSurvey}. PUBG 스쿼드 매치 역시 하나의
팀이 함께 움직이며 생존하지만, 여기에 추가 조건이 있다. 팀이 움직이는 공간 자체가
페이즈마다 축소된다. 여기서 이를 생물학적 habitat으로 주장하려는 것은 아니다. 중요한
것은 플레이 가능 영역이 줄어들고, 위치의 위험도가 페이즈에 따라 바뀌며, 남은
팀들이 점점 좁은 공간에 모인다는 구조다 \cite{fahrig2003fragmentation,nathan2008movement}.

\subsection{페이즈별 생존 모델링}

마지막으로 생존분석은 이 문제를 최종 라벨 분류가 아니라 구간별 위험으로 읽게 해
준다. 이산 시간 사건사 모델은 각 구간에서 누가 아직 위험집합에 남아 있고, 그
구간에 사건이 발생했는지를 정의한다
\cite{allison1982discrete,singer1993discrete}. MATCH-Net과 DeepHazard 같은 동적
생존 모델도 시간에 따라 바뀌는 공변량이 위험 예측에 중요함을 보였다
\cite{jarrett2018matchnet,rava2020deephazard}. 본 연구의 다음 페이즈 탈락 위험은
이 틀을 PUBG 텔레메트리에 맞게 옮긴 것이다. 차이는 위험집합의 주체가 환자나
사용자 행이 아니라, 페이즈 종료 시점까지 살아남은 스쿼드 팀이라는 점이다.
